% =============================================================================
%  SAT and Groebner Bases --- Brian Tenneson --- Version 2, July 27, 2026
%  Supersedes "Solving SAT in polynomial time?", April 11, 2023
% =============================================================================
\documentclass[11pt]{article}
\usepackage[letterpaper,margin=1in,top=1.1in,bottom=1.1in]{geometry}
\usepackage[utf8]{inputenc}
\usepackage[T1]{fontenc}
\usepackage{lmodern}
\usepackage{amsmath,amssymb,amsthm}
\usepackage{mathtools}
\usepackage{enumitem}
\usepackage{fancyhdr}
\usepackage{microtype}
\usepackage[colorlinks=true,linkcolor=blue,citecolor=blue,urlcolor=blue]{hyperref}
\usepackage{url}

\theoremstyle{plain}
\newtheorem{theorem}{Theorem}[section]
\newtheorem{proposition}[theorem]{Proposition}
\newtheorem{corollary}[theorem]{Corollary}
\theoremstyle{definition}
\newtheorem{definition}[theorem]{Definition}
\newtheorem{algorithm}[theorem]{Algorithm}
\theoremstyle{remark}
\newtheorem{remark}[theorem]{Remark}

\pagestyle{fancy}\fancyhf{}
\lhead{\footnotesize SAT and Gr\"obner Bases}
\rhead{\footnotesize Tenneson, Version 2, July 27, 2026}
\lfoot{\footnotesize Copyright \textcopyright\ 2026 Brian Tenneson. All Rights Reserved.}
\rfoot{\footnotesize btenneson2301.substack.com \quad \thepage}
\renewcommand{\headrulewidth}{0.4pt}\renewcommand{\footrulewidth}{0.4pt}
\setlength{\headheight}{14pt}

\newcommand{\FF}{\mathbf{F}_{2}}
\newcommand{\PC}{F_{\mathrm{PC}}}
\newcommand{\Res}{F_{\mathrm{res}}}

\title{\vspace{-1.2cm}\textbf{\Large SAT and Gr\"obner Bases}\\[0.5em]
{\large What the Algebraic Method Does and Does Not Buy}\\[0.6em]{\large Version 2}}
\author{Brian Tenneson\\ M.S., Applied Data Science\\ Currently Unaffiliated\\
email: \href{mailto:btenneson2301@baypath.edu}{btenneson2301@baypath.edu}}
\date{July 27, 2026}

\begin{document}
\maketitle
\thispagestyle{empty}

\begin{center}
\fbox{\parbox{0.86\textwidth}{\centering\small
Version 2 --- July 27, 2026\\[0.3em]
Supersedes \emph{Solving SAT in Polynomial Time?}, Version 1, April 11, 2023\\[0.3em]
Copyright \textcopyright\ 2026 Brian Tenneson. All Rights Reserved.\\[0.3em]
\url{https://btenneson2301.substack.com}}}
\end{center}

\vspace{0.8em}
\begin{abstract}
\noindent
Version 1 of this note proposed deciding satisfiability by arithmetizing a
propositional formula $\varphi$ as a polynomial $\tau(\varphi)$ over $\FF$,
forming the ideal $J=\langle\tau(\varphi)+1,\,x_{i}^{2}+x_{i}\rangle$, computing
a Gr\"obner basis $G$ for $J$, and testing whether $1$ reduces to $0$ modulo
$G$. It suggested that the Gr\"obner basis step might be performed in
polynomial time by the \textsc{fglm} algorithm, and hence that
$\mathrm{P}=\mathrm{NP}$.

The arithmetization and the ideal-membership test are correct, and are proved
correct here. The complexity claim is not, and this version withdraws it. Two
independent errors are identified: \textsc{fglm} converts a Gr\"obner basis
from one monomial order to another and cannot produce the first basis, which is
the expensive object; and its running time is polynomial in the dimension of
the quotient algebra $\FF[x]/J$, which counts satisfying assignments and is
exponential in the number of variables. Beyond these, the approach is closed by
a theorem rather than left open: Gr\"obner basis refutation of $J$ is exactly
proof search in the Polynomial Calculus proof system, for which exponential
lower bounds are known unconditionally.

What survives is worth stating. Polynomial Calculus polynomially simulates
resolution, so nothing is lost by working algebraically; it refutes Tseitin
formulas in polynomial size where resolution requires exponential size, so
something is strictly gained on parity-structured instances; and degree-$d$
refutations are findable in time $n^{O(d)}$, so the obstruction is degree
rather than the algebraic setting. The correct research question is not whether
this method decides \textsc{sat} in polynomial time --- it does not --- but
which families of formulas admit low-degree algebraic refutations.
\end{abstract}

\medskip
\noindent\textbf{Keywords:} satisfiability, Gr\"obner basis, Boolean ideal,
arithmetization, Polynomial Calculus, Nullstellensatz, proof complexity,
resolution, Tseitin formulas, \textsc{fglm}, degree lower bounds.

\section{The construction, and what is right about it}

Let $\varphi$ be a propositional formula over $p_{1},\dots,p_{n}$.

\begin{definition}[Arithmetization]\label{def:tau}
Define $\tau(\varphi)\in\FF[x_{1},\dots,x_{n}]$ by recursion:
\[
\tau(p_{i}):=x_{i},\qquad \tau(\neg\varphi_{1}):=1+\tau(\varphi_{1}),
\]
\[
\tau(\varphi_{1}\wedge\varphi_{2}):=\tau(\varphi_{1})\tau(\varphi_{2}),
\qquad
\tau(\varphi_{1}\vee\varphi_{2}):=1+\bigl(1+\tau(\varphi_{1})\bigr)\bigl(1+\tau(\varphi_{2})\bigr),
\]
all arithmetic in $\FF$. Put
$J(\varphi):=\langle\,\tau(\varphi)+1,\;x_{1}^{2}+x_{1},\dots,x_{n}^{2}+x_{n}\,\rangle$.
\end{definition}

\begin{proposition}[Correctness]\label{prop:correct}
For $a\in\{0,1\}^{n}$ we have $\tau(\varphi)(a)=1$ if and only if
$a\models\varphi$. Consequently
\[
\varphi\ \text{is unsatisfiable}
\iff 1\in J(\varphi)
\iff \text{the reduced Gr\"obner basis of }J(\varphi)\text{ is }\{1\}.
\]
\end{proposition}

\begin{proof}
Induction on $\varphi$. A variable is its own indicator. For negation, $1+t=1$
iff $t=0$. For conjunction, a product of elements of $\{0,1\}$ is $1$ iff both
are $1$. For disjunction, $1+(1+t_{1})(1+t_{2})=1$ iff $(1+t_{1})(1+t_{2})=0$
iff some $t_{i}=1$.

The field equations $x_{i}^{2}+x_{i}$ force every common zero into
$\{0,1\}^{n}$, so the variety of $J(\varphi)$ is the set of satisfying
assignments. That variety is empty iff $J(\varphi)$ is the unit ideal --- the
ideal is zero-dimensional and radical, since it contains the field equations,
so the weak Nullstellensatz applies without a radical correction --- and an
ideal is the unit ideal iff its reduced Gr\"obner basis in any monomial order
is $\{1\}$.
\end{proof}

\begin{algorithm}[The procedure of Version 1]\label{alg:v1}
On input $\varphi$: compute $\tau(\varphi)$; form $J(\varphi)$; compute a
Gr\"obner basis $G$ of $J(\varphi)$; reduce $1$ modulo $G$; report
\emph{unsatisfiable} if the reduction is $0$ and \emph{satisfiable} otherwise.
\end{algorithm}

\begin{remark}
Algorithm~\ref{alg:v1} is a correct decision procedure, by
Proposition~\ref{prop:correct}. Everything at issue is its running time, and
that is entirely concentrated in the third step.
\end{remark}

\section{Why the complexity claim fails}

Version 1 proposed to perform the Gr\"obner basis computation ``in polynomial
time using the \textsc{fglm} algorithm.'' There are two separate problems with
this, and then a third that is fatal independently of the first two.

\subsection{\textsc{fglm} is a change of order, not a computation from scratch}

The algorithm of Faug\`ere, Gianni, Lazard and Mora \cite{fglm} takes as
\emph{input} a Gr\"obner basis of a zero-dimensional ideal with respect to one
monomial order and returns a Gr\"obner basis with respect to another --- in
practice, from a fast degree order to the lexicographic order convenient for
elimination. It presupposes the object whose cost is in question. Using it as
the step that computes $G$ is circular: the first basis must be obtained by
Buchberger's algorithm or a descendant such as $F_{4}$ or $F_{5}$, and that is
where the work is.

\subsection{Its running time is polynomial in the wrong parameter}

\textsc{fglm} runs in time $O(nD^{3})$, where
\[
D=\dim_{\FF}\bigl(\FF[x_{1},\dots,x_{n}]/J\bigr).
\]
For a zero-dimensional radical ideal $D$ is the number of points of the
variety, so here $D$ is the number of satisfying assignments of $\varphi$.
This is at most $2^{n}$ and is routinely close to it. Being polynomial in $D$
is not being polynomial in $n$: already the field equations alone give
$\dim_{\FF}\FF[x]/\langle x_{i}^{2}+x_{i}\rangle=2^{n}$.

The mistake is a substitution of parameters. ``Polynomial time'' in the
computer-algebra literature is frequently relative to the quotient dimension or
to the B\'ezout number, not to the bit length of the input, and the two
coincide only for ideals with few solutions.

\subsection{The approach is closed by a theorem}

The decisive point is not that the two objections above might be circumvented
by a better algorithm. It is that the method as a whole is a known proof system
with known lower bounds.

\begin{definition}[Polynomial Calculus]\label{def:pc}
Over a field $K$, the \emph{Polynomial Calculus} $\PC$ has as its formulas the
multilinear polynomials of $K[x_{1},\dots,x_{n}]/\langle x_{i}^{2}-x_{i}\rangle$
and as its rules
\[
(f,g)\longmapsto \alpha f+\beta g\ \ (\alpha,\beta\in K),
\qquad
f\longmapsto x_{i}f .
\]
A \emph{refutation} of a set of polynomials is a derivation of $1$. Its
\emph{size} is the total number of monomials and its \emph{degree} is the
largest total degree occurring in it.
\end{definition}

\begin{proposition}\label{prop:gbispc}
Every run of Buchberger's algorithm on $J(\varphi)$ that outputs $\{1\}$ is a
$\PC$ refutation of the generators of $J(\varphi)$, of size at most the number
of monomials the run writes down.
\end{proposition}

\begin{proof}
The algorithm forms $S$-polynomials and reduces them. An $S$-polynomial of $f$
and $g$ is $m_{f}f-m_{g}g$ for monomials $m_{f},m_{g}$, which is a composition
of multiplications by variables followed by one linear combination; a reduction
step subtracts a monomial multiple of a basis element, likewise. Every
intermediate polynomial produced is therefore derivable by the two rules of
Definition~\ref{def:pc} from the generators, and $1$ is among them exactly when
the output basis is $\{1\}$.
\end{proof}

\begin{theorem}[Lower bounds]\label{thm:lower}
Let $\Gamma_{n}$ be the arithmetized pigeonhole principle
$\mathrm{PHP}_{n}^{n-1}$. Every $\PC$ refutation of $\Gamma_{n}$ has degree
$\Omega(n)$, and there is $c>0$ such that every $\PC$ refutation of
$\Gamma_{n}$ has size at least $2^{cn}$. The corresponding statement holds
almost surely for random $3$-CNF formulas at constant density above the
satisfiability threshold.
\end{theorem}

\begin{proof}[Attribution]
The degree bound is Razborov's \cite{razborov}; the size bound follows by the
size--degree tradeoff of Impagliazzo, Pudl\'ak and Sgall \cite{ips}, whose
paper is titled for this consequence. The random-formula statement is
Ben-Sasson and Impagliazzo \cite{bsi}. The proof system was introduced, and its
identification with the Gr\"obner basis algorithm made, by Clegg, Edmonds and
Impagliazzo \cite{clegg}.
\end{proof}

\begin{corollary}\label{cor:nopoly}
Algorithm~\ref{alg:v1} does not run in time polynomial in $|\varphi|$, and no
variant of it that decides unsatisfiability by exhibiting $1$ in the ideal
does. This is unconditional: it assumes no complexity-theoretic hypothesis, and
in particular does not depend on $\mathrm{P}\ne\mathrm{NP}$.
\end{corollary}

\begin{proof}
On input $\Gamma_{n}$ the algorithm's run is a $\PC$ refutation by
Proposition~\ref{prop:gbispc}, and its running time is at least the number of
monomials it writes. By Theorem~\ref{thm:lower} that number is $2^{\Omega(n)}$,
while $|\Gamma_{n}|$ is polynomial in $n$.
\end{proof}

\begin{remark}[The claim of Version 1 is withdrawn]
Corollary~\ref{cor:nopoly} settles the question the earlier version left open,
in the negative. The route from Gr\"obner bases to $\mathrm{P}=\mathrm{NP}$ is
not merely unproved; it is blocked by a theorem proved in 1999. It should be
added that this is a good outcome for the method as mathematics: a technique
with an exponential lower bound is a technique whose behaviour is understood.
\end{remark}

\section{What survives}

The lower bounds are worst-case statements, and the algebraic method retains
three properties that clause-based methods do not have. These, and not the
complexity claim, are the reason to be interested in it.

\begin{theorem}[Simulation]\label{thm:sim}
$\PC$ polynomially simulates resolution over every field, and the degree of the
simulating refutation equals the width of the resolution refutation simulated.
Hence nothing is given up by working algebraically.
\end{theorem}

\begin{proof}[Attribution and sketch]
Translate a clause $\bigvee_{i\in P}p_{i}\vee\bigvee_{j\in N}\neg p_{j}$ to the
monomial $\prod_{i\in P}(1-x_{i})\prod_{j\in N}x_{j}$, which vanishes exactly
on the assignments satisfying the clause. A resolution step becomes a bounded
number of applications of the two rules. Due to Clegg, Edmonds and Impagliazzo
\cite{clegg}.
\end{proof}

\begin{theorem}[Strict gain on parity structure]\label{thm:tseitin}
Let $T_{n}$ be the Tseitin formula of a bounded-degree expander graph on $n$
vertices with odd total charge. Then $T_{n}$ has $\PC$ refutations of
polynomial size over $\FF$, while every resolution refutation of $T_{n}$ has
size $2^{\Omega(n)}$.
\end{theorem}

\begin{proof}[Attribution and sketch]
Over $\FF$ the parity constraint at a vertex $v$ is the linear polynomial
$\sum_{e\ni v}x_{e}+\sigma(v)$. Summing these over all vertices makes each edge
variable occur twice and cancel, leaving $\sum_{v}\sigma(v)=1$: a refutation
using one linear combination per vertex and no multiplications whatever. The
exponential resolution lower bound is Urquhart's \cite{urquhart}.
\end{proof}

\begin{corollary}
$\Res\preceq_{\mathrm{p}}\PC$ but $\PC\not\preceq_{\mathrm{p}}\Res$: the
algebraic system is strictly stronger than resolution.
\end{corollary}

\begin{theorem}[Bounded degree is tractable]\label{thm:degree}
For fixed $d$ there is an algorithm running in time $n^{O(d)}$ which finds a
$\PC$ refutation of degree at most $d$ whenever one exists.
\end{theorem}

\begin{proof}[Sketch]
Multilinear polynomials of degree $\le d$ in $n$ variables span a space of
dimension $\sum_{j\le d}\binom{n}{j}=n^{O(d)}$. Closing the input under
multiplication by variables and taking linear spans, discarding anything of
degree above $d$, is Gaussian elimination on a matrix of that side length; $1$
is derivable in degree $d$ iff it lands in the span \cite{clegg}.
\end{proof}

\begin{remark}[The question that replaces the original one]\label{rem:newq}
Theorem~\ref{thm:degree} localizes the difficulty. The algebraic setting is not
itself the obstruction; \emph{degree} is. Where a family of formulas admits
refutations of bounded degree, those refutations are found in polynomial time
for each fixed bound, and Theorem~\ref{thm:tseitin} exhibits a natural family
admitting refutations of degree $1$. The productive question is therefore:

\begin{center}
\emph{Which families of propositional formulas admit low-degree Polynomial
Calculus refutations?}
\end{center}

\noindent This is answerable, is being answered, and has the property the
original question lacked --- that partial progress on it is progress. By
contrast, Theorem~\ref{thm:lower} tells us in advance that the answer is not
``all of them.''
\end{remark}

\begin{remark}[On the Boolean-ring presentation]
Carrying the field equations as explicit generators, as in
Definition~\ref{def:tau}, is one of several presentations. Boolean Gr\"obner
basis theory works directly in the quotient ring, and admits coefficient rings
that are themselves Boolean rather than fields \cite{sato}. This is often
substantially better in practice and is the right setting for implementation.
It does not affect any statement above: the lower bounds of
Theorem~\ref{thm:lower} are proved for the proof system, not for a particular
data structure.
\end{remark}

\section{Summary}

\begin{center}
\begin{tabular}{@{}p{0.42\textwidth}p{0.5\textwidth}@{}}
\hline
\textbf{Claim of Version 1} & \textbf{Status} \\
\hline
The arithmetization $\tau$ is correct
& Correct; Proposition~\ref{prop:correct}. \\
$\varphi$ unsatisfiable iff $1$ reduces to $0$ mod $G$
& Correct; Proposition~\ref{prop:correct}. \\
\textsc{fglm} computes $G$ in polynomial time
& False twice over: \textsc{fglm} converts an existing basis, and is polynomial in the quotient dimension, which is exponential in $n$. \\
Hence $\mathrm{P}=\mathrm{NP}$
& Withdrawn. Corollary~\ref{cor:nopoly} shows the method is exponential on the pigeonhole family, unconditionally. \\
\hline
\textbf{Retained} & \\
\hline
Algebraic methods subsume clause methods
& Theorem~\ref{thm:sim}. \\
Algebraic methods strictly beat resolution somewhere
& Theorem~\ref{thm:tseitin}, on Tseitin formulas. \\
Low-degree refutations are efficiently findable
& Theorem~\ref{thm:degree}. \\
\hline
\end{tabular}
\end{center}

\bigskip
\noindent This note is a companion to Section 10 of \emph{Depths of a
Simulation and Formalized Self-Awareness}, Merged Edition, where the same
material is developed inside the formal-system vocabulary $F=(x,y,z)$ and the
Tseitin separation is recorded as a second unconditional instance of the
Two-System Size-Speedup Principle.

\section*{Acknowledgements}
The author thanks AI assistants used as editorial and verification tools during
the preparation of this revision, in particular for the identification of the
two \textsc{fglm} errors and the location of the relevant lower bounds. All
mathematical claims, interpretations, and final decisions remain the
responsibility of the author.

\begin{thebibliography}{99}
\bibitem{bsi} Ben-Sasson, E., \& Impagliazzo, R. (1999). Random CNF's are hard for the polynomial calculus. In \textit{Proceedings of the 40th Annual Symposium on Foundations of Computer Science} (pp. 415--421). IEEE.
\bibitem{clegg} Clegg, M., Edmonds, J., \& Impagliazzo, R. (1996). Using the Groebner basis algorithm to find proofs of unsatisfiability. In \textit{Proceedings of the Twenty-Eighth Annual ACM Symposium on the Theory of Computing} (pp. 174--183). ACM.
\bibitem{cr} Cook, S. A., \& Reckhow, R. A. (1979). The relative efficiency of propositional proof systems. \textit{Journal of Symbolic Logic}, 44(1), 36--50.
\bibitem{fglm} Faug\`ere, J.-C., Gianni, P., Lazard, D., \& Mora, T. (1993). Efficient computation of zero-dimensional Gr\"obner bases by change of ordering. \textit{Journal of Symbolic Computation}, 16(4), 329--344.
\bibitem{ips} Impagliazzo, R., Pudl\'ak, P., \& Sgall, J. (1999). Lower bounds for the polynomial calculus and the Gr\"obner basis algorithm. \textit{Computational Complexity}, 8(2), 127--144.
\bibitem{razborov} Razborov, A. A. (1998). Lower bounds for the polynomial calculus. \textit{Computational Complexity}, 7(4), 291--324.
\bibitem{sato} Sato, Y., Inoue, S., Suzuki, A., Nabeshima, K., \& Sakai, K. (2011). Boolean Gr\"obner bases. \textit{Journal of Symbolic Computation}, 46(5), 622--632.
\bibitem{urquhart} Urquhart, A. (1987). Hard examples for resolution. \textit{Journal of the ACM}, 34(1), 209--219.
\end{thebibliography}

\end{document}
